% ============================================================
% 第 14 章 三相感应电动机的运行特性（教材第四篇，印刷页 150--166）
% 节名：第一节 三相感应电动机转矩—转差率曲线；第二节 三相感应电动机的机械特性和工作特性；
%       第三节 深槽和双笼感应电动机；第四节 三相感应电动机的启动及电磁制动；第五节 三相感应电动机的调速
% ============================================================
\chapter{三相感应电动机的运行特性}

第 13 章给出了电磁转矩的表达式。本章把它画成曲线（转矩—转差率曲线、机械特性），再由曲线解释三个工程问题：\textbf{能不能稳定运行}、\textbf{怎么起动}、\textbf{怎么调速}。这三个问题也是感应电动机全部工程应用的落脚点。

\section{三相感应电动机转矩—转差率曲线}

在额定电压和额定频率下，电磁转矩与转差率的关系 $T=f(s)$ 称为\textbf{转矩—转差率曲线}。由第 13 章
\begin{equation}
T=\frac{m_1pU_1^{2}\dfrac{R_2'}{s}}{\Omega_1\left[\left(R_1+\dfrac{R_2'}{s}\right)^{2}+\left(X_{1\sigma}+X_{2\sigma}'\right)^{2}\right]} .
\end{equation}

\subsection{曲线的三个特征点}
\begin{enumerate}
  \item \textbf{同步点}（$s=0$，$n=n_1$）：$T=0$。转子与旋转磁场无相对运动，无感应电动势、无电流、无转矩。
  \item \textbf{最大转矩点}（$s=s_m$）：$T=T_{\max}$。对 $s$ 求导并令其为零，得\textbf{临界转差率}
  \begin{equation}
  s_m=\frac{R_2'}{\sqrt{R_1^{2}+\left(X_{1\sigma}+X_{2\sigma}'\right)^{2}}} ,
  \end{equation}
  代回转矩式得\textbf{最大转矩}
  \begin{equation}
  T_{\max}=\frac{m_1pU_1^{2}}{2\Omega_1\left[R_1+\sqrt{R_1^{2}+\left(X_{1\sigma}+X_{2\sigma}'\right)^{2}}\right]} .
  \end{equation}
  \item \textbf{起动点}（$s=1$，$n=0$）：$T=T_{st}$ 称为\textbf{起动转矩}，是曲线与 $s=1$ 的交点。
\end{enumerate}

\keyformula{s_m=\frac{R_2'}{\sqrt{R_1^{2}+X_k^{2}}},\qquad
T_{\max}=\frac{m_1pU_1^{2}}{2\Omega_1\left[R_1+\sqrt{R_1^{2}+X_k^{2}}\right]},\qquad X_k=X_{1\sigma}+X_{2\sigma}'}

\subsection{两个重要比例关系}
\begin{itemize}
  \item $T_{\max}\propto U_1^{2}$：\textbf{最大转矩与电压平方成正比}。电网电压降低 10\%，最大转矩降低约 19\%，这是电动机"低电压时容易堵转烧毁"的原因。
  \item $s_m\propto R_2'$：\textbf{临界转差率与转子电阻成正比}，而 $T_{\max}$ \textbf{与 $R_2'$ 无关}（忽略 $R_1$ 时）。这一条是绕线型电动机"转子串电阻改善起动性能"的理论依据：增大 $R_2'$ 可使 $s_m$ 向 1 靠近，从而提高起动转矩，而最大转矩不变。
\end{itemize}

\subsection{稳定运行区}
把曲线分为两段：
\begin{itemize}
  \item $0<s<s_m$（转速较高段）：$T$ 随 $s$ 增大而增大，$\dfrac{\mathrm{d}T}{\mathrm{d}n}<0$，为\textbf{稳定运行区}。负载增加时转速略降、转矩增大，自动达到新的平衡。
  \item $s_m<s<1$（转速较低段）：$T$ 随 $s$ 增大反而减小，为\textbf{不稳定区}。若负载转矩超过 $T_{\max}$，电机将堵转。
\end{itemize}
稳定运行的判据是：负载扰动消失后能回到原平衡点，即 $\dfrac{\mathrm{d}T}{\mathrm{d}n}<\dfrac{\mathrm{d}T_L}{\mathrm{d}n}$。额定运行点一般在 $s_N=0.01\sim0.05$，距 $s_m$ 尚有裕量，故过载能力
\begin{equation}
\lambda=\frac{T_{\max}}{T_N}=1.6\sim2.5 .
\end{equation}

\section{三相感应电动机的机械特性和工作特性}

\subsection{机械特性}
把 $T=f(s)$ 换成 $n=f(T)$ 即得\textbf{机械特性}。它近似为一条略微下倾的直线（工作段），在 $T_{\max}$ 处出现拐点，之后转速迅速下降到零（堵转）。

机械特性是分析电动机与负载配合的工具：把电动机的 $n=f(T)$ 与负载的 $n=f(T_L)$ 画在同一图上，交点即稳态工作点。

\subsection{工作特性}
工作特性指 $U_1=U_N$、$f_1=f_N$ 时，转速 $n$、电磁转矩 $T$、定子电流 $I_1$、功率因数 $\cos\varphi_1$、效率 $\eta$ 随输出功率 $P_2$ 变化的曲线。

\begin{itemize}
  \item \textbf{转速特性} $n=f(P_2)$：从空载到额定负载，转速下降很少（$s$ 由 0 增到 $s_N$），曲线略向下倾斜，属"硬特性"。
  \item \textbf{转矩特性} $T=f(P_2)$：$T=T_0+P_2/\Omega$，近似为一条略上翘的直线。
  \item \textbf{定子电流特性} $I_1=f(P_2)$：空载时 $I_1=I_0$（较大，占额定 20\%～50\%）；随负载增加，$I_1$ 近似按线性增长。
  \item \textbf{功率因数特性} $\cos\varphi_1=f(P_2)$：空载时很低（0.1～0.2，因励磁电流大），随负载增加迅速上升，在额定负载附近达最大（0.8～0.9），过载后又下降。
  \item \textbf{效率特性} $\eta=f(P_2)$：空载时为零，随负载增加而上升，在可变损耗（铜耗）等于不变损耗（铁耗+机械损耗）时达到最大，一般在 $(0.7\sim0.8)P_N$ 附近。
\end{itemize}

\begin{center}
\begin{tabular}{@{}lll@{}}
\toprule
特性 & 变化规律 & 原因 \\
\midrule
$n=f(P_2)$ & 略下降（硬特性） & $s_N$ 仅 1\%～5\% \\
$T=f(P_2)$ & 近似直线略上翘 & $T=T_0+P_2/\Omega$ \\
$I_1=f(P_2)$ & 空载较大，随负载近线性增大 & 定子电流含励磁分量与负载分量 \\
$\cos\varphi_1=f(P_2)$ & 空载低、额定附近最高 & 励磁电流占比随负载下降 \\
$\eta=f(P_2)$ & 先升后降，$(0.7\sim0.8)P_N$ 最高 & 可变损耗 $=$ 不变损耗 \\
\bottomrule
\end{tabular}
\end{center}

\corepoint{感应电动机"轻载运行不经济"的根子在功率因数：空载时 $I_1$ 中绝大部分是无功的励磁电流，$\cos\varphi_1$ 只有 0.1～0.2；负载增加后有功分量上升，$\cos\varphi_1$ 才提高。所以"大马拉小车"不仅效率低，还使电网功率因数变差——这是工程上按负载选配电机容量的理由。}

\section{深槽和双笼感应电动机}

普通笼型电动机起动时 $s=1$，$X_{2\sigma}$ 相对较大、转子功率因数低，起动转矩不大而起动电流很大（$I_{st}=(4\sim7)I_N$）。改善起动性能的关键是\textbf{起动时增大转子电阻、运行时恢复小电阻}——利用\textbf{集肤效应}可以自动实现这一点。

\subsection{集肤效应}
当转子导条中电流频率较高时，交流电流因电感作用趋向于导条\textbf{上部（槽口侧）}流动，有效截面减小、等效电阻增大、漏抗减小，这种现象称为\textbf{集肤效应}（挤流效应）。集肤效应的强弱与频率有关：频率越高越显著，频率越低越不明显。

而转子频率 $f_2=sf_1$：\textbf{起动瞬间} $s=1$，$f_2=f_1=50$\,Hz，集肤效应强，转子电阻大 $\Rightarrow$ 起动转矩大、起动电流小；\textbf{正常运行} $s$ 很小，$f_2$ 仅 1～3\,Hz，集肤效应几乎消失，转子电阻回到小值 $\Rightarrow$ 效率高、转差率小。这就是深槽与双笼转子的巧妙之处：\textbf{不需要外部控制，靠频率自动切换}。

\subsection{深槽转子}
转子槽形做得又深又窄（高宽比可达 10～12），导条为高而窄的矩形。起动时集肤效应使电流集中在上部，等效电阻大；运行时电流均匀分布，电阻小。

\subsection{双笼转子}
转子上有两套笼：\textbf{上笼}（起动笼）用黄铜或铝青铜等电阻率高的材料，截面小、靠近槽口，漏抗小；\textbf{下笼}（运行笼）用紫铜，截面大、埋在槽底，漏抗大。
\begin{itemize}
  \item 起动时 $f_2$ 高，下笼漏抗大、电流被挤到上笼，上笼电阻大 $\Rightarrow$ \textbf{起动转矩大}；
  \item 运行时 $f_2$ 低，两笼电流按电阻反比分配，下笼电阻小、承担大部分电流 $\Rightarrow$ \textbf{运行效率高}。
\end{itemize}
双笼电动机的转矩—转差率曲线呈现"两段上凸"的形状，兼顾了起动与运行性能。

\section{三相感应电动机的启动及电磁制动}

\subsection{对起动性能的要求}
\begin{enumerate}
  \item \textbf{起动转矩足够大}（$T_{st}/T_N$ 一般要求 1.0～2.0），能带动负载起转；
  \item \textbf{起动电流不要太大}（普通笼型 $I_{st}/I_N=4\sim7$），以免电网电压跌落、绕组过热；
  \item 起动设备简单、经济、可靠。
\end{enumerate}
起动瞬间 $s=1$、转子支路阻抗最小，故电流最大而转矩并不最大——这是感应电动机起动问题的根源。

\subsection{笼型电动机的起动方法}

\subsubsection{直接起动（全压起动）}
把定子直接接入额定电压。方法最简单、起动转矩最大，但起动电流大。适用于小容量电动机（一般满足"电动机容量不超过电源变压器容量的 20\%～30\%"这类经验条件）。

\subsubsection{星—三角（Y—$\Delta$）降压起动}
运行时为 $\Delta$ 接的电动机，起动时先接成 Y，起动完毕再切换为 $\Delta$。由于 Y 接时每相电压降为 $1/\sqrt3$，故
\begin{equation}
\frac{I_{stY}}{I_{st\Delta}}=\frac13,\qquad
\frac{T_{stY}}{T_{st\Delta}}=\left(\frac{1}{\sqrt3}\right)^{2}=\frac13 .
\end{equation}
即\textbf{起动电流和起动转矩都降为直接起动的 1/3}。设备简单、成本低，但只适用于正常运行时为 $\Delta$ 接的电机，且起动转矩降低较多，只宜轻载起动。

\subsubsection{自耦变压器降压起动}
用自耦变压器把电压降低后加到定子，起动完毕再切换到全压。若抽头电压比为 $k_A<1$，则
\begin{equation}
\frac{I_{st}}{I_{st,\text{直接}}}=k_A^{2},\qquad
\frac{T_{st}}{T_{st,\text{直接}}}=k_A^{2} .
\end{equation}
有 55\%、64\%、73\% 等抽头可选，比 Y—$\Delta$ 灵活，可用于 Y 接或 $\Delta$ 接的电机；缺点是设备体积大、价格高。

\subsubsection{软起动与变频起动}
用晶闸管调压（软起动器）使电压平滑上升，或用变频器从低频低压起动，可把起动电流限制在 1.5～2 倍额定值以内，且起动平滑、无冲击，是现代主流方案。

\subsection{绕线型电动机的起动}
绕线型转子可外串电阻。由 $s_m\propto R_2'$、$T_{\max}$ 不变可知：串入适当电阻既\textbf{减小起动电流}，又\textbf{提高起动转矩}（可做到 $T_{st}=T_{\max}$）。起动过程中逐段切除电阻，可使转矩保持在较大范围内，是起重、轧钢等重载起动的理想方案。

\subsection{电磁制动}
\begin{enumerate}
  \item \textbf{反接制动}：把定子任意两相对调，旋转磁场反向，转差率 $s>1$，转矩变为制动转矩，转速迅速下降；到零附近必须及时切断电源，否则电机会反转。制动强烈但能量全部消耗在转子回路。
  \item \textbf{能耗制动}：切断三相交流电源，立即在定子绕组中通入直流电流，产生静止磁场，转子切割该磁场产生制动转矩。制动平稳、无冲击，能量也变为热。
  \item \textbf{回馈制动}：电机转速超过同步转速（$s<0$），进入发电机状态，把机械能回馈电网。多用于起重机下放重物、变频调速减速过程，最经济。
\end{enumerate}

\section{三相感应电动机的调速}

由 $n=(1-s)n_1=(1-s)\dfrac{60f_1}{p}$ 可知，调速只有三条路径：\textbf{改变极对数 $p$}、\textbf{改变频率 $f_1$}、\textbf{改变转差率 $s$}。

\subsection{变极调速}
改变定子绕组的接法（如 Y—YY、$\Delta$—YY）使极对数成倍变化，同步转速随之成倍变化，故为\textbf{有级调速}。优点是不需额外设备、效率高；缺点是级差大、不能连续调节，多用于机床、风机等。

\subsection{变频调速}
改变供电频率 $f_1$ 即可连续改变 $n_1$，是目前性能最好的调速方式。关键约束是\textbf{保持磁通恒定}：由 $U_1\approx E_1=4.44f_1N_1k_{w1}\Phi_m$ 可知，若只降频不降压，$\Phi_m$ 将增大导致磁路饱和、励磁电流剧增，故需按
\begin{equation}
\frac{U_1}{f_1}=\text{常数}
\end{equation}
进行"恒磁通"控制（基频以下恒转矩调速）；基频以上因电压不能超过额定值，转为"恒功率"调速（弱磁升速）。变频调速范围宽、平滑、效率高，是现代交流调速的主流。

\subsection{改变转差率调速}
\begin{itemize}
  \item \textbf{转子串电阻调速}（绕线型）：串电阻使 $s_m$ 增大，同一转矩下转差率增大、转速下降。设备简单，但转差功率消耗在电阻上，效率随转速降低而下降。
  \item \textbf{串级调速}（绕线型）：把转子回路的转差功率经整流、逆变后回馈电网或用于驱动，效率高，但设备复杂。
  \item \textbf{改变定子电压调速}：降低 $U_1$ 使 $T_{\max}\propto U_1^2$ 下降，工作点沿负载特性下移，转速降低。方法简单但调速范围小、低速时损耗大，多用于风机类平方转矩负载。
\end{itemize}

\begin{center}
\begin{tabular}{@{}p{2.3cm}p{3.4cm}p{3.8cm}p{3.4cm}@{}}
\toprule
调速方式 & 依据 & 特点 & 适用 \\
\midrule
变极 & 改变 $p$，$n_1\propto1/p$ & 有级、高效、便宜 & 机床、风机 \\
变频 & 改变 $f_1$，$n_1\propto f_1$ & 无级、范围宽、效率高 & 现代主流，需变频器 \\
转子串电阻 & 改变 $s$（$s_m\propto R_2'$） & 简单，转差功率损耗大 & 绕线型，短时调速 \\
串级 & 转差功率回馈 & 效率高、设备复杂 & 绕线型大容量 \\
改变定子电压 & $T_{\max}\propto U_1^2$ & 范围小、低速损耗大 & 风机、泵类 \\
\bottomrule
\end{tabular}
\end{center}

\corepoint{本章的两条主线要抓牢：① \textbf{转矩公式的三个比例}——$T\propto U_1^2$、$T_{\max}\propto U_1^2$ 且与 $R_2'$ 无关、$s_m\propto R_2'$，几乎所有起动与调速问题都由这三条推出；② \textbf{起动与调速都在"动 $s$ 或动 $n_1$"}——降压、串电阻是动 $s$（并牺牲转矩或效率），变极、变频是动 $n_1$（不改转差率，故效率高）。答题时先说清"我动的是 $n_1$ 还是 $s$"，思路就不会乱。}
